\section*{Capítulo \ref{ch:g4:magnets-and-compass} --- Los imanes y la brújula}

\begin{solution}{exo:g4:magnets-and-compass:1}
Los \omterm{def:g4:magnets-and-compass:poles}{polos} son los dos extremos fuertes del \omterm{def:g1:magnets:magnet}{imán}: el
\omterm{def:g4:magnets-and-compass:poles}{polo norte} (N, a menudo rojo) y el \omterm{def:g4:magnets-and-compass:poles}{polo sur} (S). Los clips se
amontonan en los \omterm{def:g4:magnets-and-compass:poles}{polos} --- en los extremos ---, no en el centro.
\end{solution}

\begin{solution}{exo:g4:magnets-and-compass:2}
Los \omterm{def:g4:magnets-and-compass:poles}{polos} opuestos se atraen; los iguales se repelen. Un golpe
seco significa que se estaban mirando \omterm{def:g4:magnets-and-compass:poles}{polos} opuestos: una N frente a
una S.
\end{solution}

\begin{solution}{exo:g4:magnets-and-compass:3}
N hacia N: se apartan (los \omterm{def:g4:magnets-and-compass:poles}{polos} iguales se repelen). N hacia S: se
atraen y se pegan de golpe.
\end{solution}

\begin{solution}{exo:g4:magnets-and-compass:4}
Porque un \omterm{def:g1:magnets:magnet}{imán} gigante está poniendo sus \omterm{def:g4:magnets-and-compass:poles}{polos} en fila: la propia
Tierra es un \omterm{def:g1:magnets:magnet}{imán}, con \omterm{def:g4:magnets-and-compass:poles}{polos} magnéticos cerca de la parte de
arriba y de la de abajo del globo.
\end{solution}

\begin{solution}{exo:g4:magnets-and-compass:5}
Norte, este, sur y oeste. Mirando al norte, el este queda a tu
derecha.
\end{solution}

\begin{solution}{exo:g4:magnets-and-compass:6}
La aguja es un \omterm{def:g1:magnets:magnet}{imán} diminuto y obedece al tirón fuerte \emph{más
cercano}. Un \omterm{def:g1:magnets:magnet}{imán} o una gran \omterm{def:g3:measuring-mass:mass}{masa} de hierro que estén cerca ganan
al lejano \omterm{def:g1:magnets:magnet}{imán} de la Tierra, y la aguja miente sobre dónde está el
norte.
\end{solution}

\begin{solution}{exo:g4:magnets-and-compass:7}
Se frota la aguja con un \omterm{def:g4:magnets-and-compass:poles}{polo} de un \omterm{def:g1:magnets:magnet}{imán} veinte veces en el
mismo sentido, apartando el \omterm{def:g1:magnets:magnet}{imán} entre pasada y pasada --- así se
convierte en un imancito. Se pone sobre una rodaja de corcho que flote
en un cuenco de agua: el corcho flotante es un \omterm{def:g3:levers-and-scales:lever}{punto de apoyo} hecho
de agua, y la aguja gira hasta apuntar norte--sur.
\end{solution}

\begin{solution}{exo:g4:magnets-and-compass:8}
Apunta al norte. Una \omterm{def:g4:magnets-and-compass:compass}{brújula} cuya aguja coincide con la \omterm{def:g1:light-and-shadows:shadow}{sombra} del
mediodía está diciendo la verdad --- y, en un día soleado, la \omterm{def:g1:light-and-shadows:shadow}{sombra}
sola puede servir de \omterm{def:g4:magnets-and-compass:compass}{brújula} aproximada.
\end{solution}

\begin{solution}{exo:g4:magnets-and-compass:9}
Funciona con niebla, bajo las nubes, en noches sin luna y lejos de
cualquier costa --- la aguja no le pide nada al cielo. Los marineros
pudieron por fin mantener el rumbo por mar abierto con cualquier
tiempo.
\end{solution}

\begin{solution}{exo:g4:magnets-and-compass:10}
La hebilla de hierro del amigo gana al \omterm{def:g1:magnets:magnet}{imán} de la Tierra a corta
distancia: la aguja apunta a la hebilla y no al norte. Si se fiara de
ella, se desviaría --- el falso ``norte'' queda hacia su izquierda,
así que torcería su camino hacia la izquierda, es decir, hacia el
oeste. El arreglo: leer la \omterm{def:g4:magnets-and-compass:compass}{brújula} lejos de la hebilla (a unos
pasos de distancia) y después caminar.
\end{solution}

\begin{solution}{exo:g4:magnets-and-compass:11}
La aguja corta que se ha frotado tiene \emph{dos} \omterm{def:g4:magnets-and-compass:poles}{polos} --- sus
extremos se comportan como una N y una S (haz \omterm{def:g1:floating-and-sinking:words}{flotar} dos agujas así
y se pegarán punta con punta, atrayéndose por los extremos opuestos).
O sea que acortar un \omterm{def:g1:magnets:magnet}{imán} no le quita un \omterm{def:g4:magnets-and-compass:poles}{polo}, y cortar uno por
la mitad solo produce dos imanes cortos completos: la caza del
\omterm{def:g4:magnets-and-compass:poles}{polo} solitario a base de cortar no acaba nunca --- cada corte te
devuelve los dos \omterm{def:g4:magnets-and-compass:poles}{polos}.
\end{solution}
